% !TEX root = ../main.tex
\section{Introducción}
\label{sec:introduccion}
\label{sec:resumen}
Los escribas digitales generan un borrador de la nota clínica a partir de la conversación entre el profesional y el paciente. Su promesa es reducir el trabajo de documentación y permitir que el profesional dedique más atención a la consulta. La nota requiere revisión y aprobación antes de incorporarse a la historia clínica~\cite{7MXFHMBA}.

La rápida expansión del mercado refleja el interés por aliviar esa carga. Los estudios describen mejoras de bienestar y, en algunos contextos, de tiempo de documentación, con resultados variables~\cite{FKHM4UFL,4PZBAAZ7}. El beneficio depende también del trabajo de revisión y del efecto sobre la relación asistencial. Una encuesta de pacientes de UC Davis recoge expectativas de mayor atención y preocupaciones sobre privacidad y exactitud. Al ser previa a la implantación, informa sobre aceptación prevista, no sobre cambios observados en la comunicación~\cite{54NFKG36}.

Este informe compara \textbf{43 proveedores} para un hospital con \textbf{Dedalus HCIS y atención en español}. Las prioridades son protección de datos, ciberseguridad, calidad, integración y coste total. ENS de categoría ALTA se mantiene como criterio prioritario, sujeto a definición institucional. La documentación aún no permite confirmar un candidato con integración ambiental operativa con HCIS y alcance ENS íntegramente verificado (sección~\ref{sec:candidatos}).

Se propone avanzar de forma progresiva: definir condiciones de uso, valorar un conjunto local de consultas de prueba y comparar entre uno y tres proveedores antes de ampliar el uso asistencial. La evaluación deberá comprobar el beneficio en español y en el flujo completo hasta HCIS, incluyendo la revisión de las notas. Servicios participantes, política de datos, presupuesto, protocolo y umbrales de aceptación permanecen abiertos.
